\documentclass[12pt]{article}
\usepackage[utf8]{inputenc}
\usepackage[T1]{fontenc}
\usepackage{hyperref}
\title{Syriac Language Resources for K-12 Curriculum Development}
\author{Compiled by Claude Sonnet 5.5 (claude-sonnet-5-5) for Andr\'as Kornai}
\date{}
\begin{document}
\maketitle

\section{Introduction}
The prompt gave only the ISO 639-3 code syr, which is the macrolanguage code for Syriac, with no glottocode, family, or speaker data. Search results concern mainly Classical Syriac, the literary and liturgical Aramaic of Middle Eastern Christians. Resources below are for Classical Syriac, as found in search. Modern spoken varieties (Neo-Aramaic) were not separately researched and may need their own dossiers.

\section{Primary Educational Resources}
\begin{itemize}
\item Steven C. Hallam, \emph{Basics of Classical Syriac}: a beginning grammar, workbook and lexicon in one volume, described as suitable for independent learners and classrooms. No URL for this title appeared in the search results.
\item George Kiraz, \emph{The New Syriac Primer}, 2nd ed. (Gorgias Press, 2013). No URL for this title appeared in the search results.
\item UCLouvain university course on Syriac (taught in French), covering morphology and grammar, the Estrangelo, Serto and East Syriac scripts, and reading of simple vocalized texts. \url{https://uclouvain.be/en-cours-2026-lglor1644a}
\end{itemize}
No school-level (K-12) Syriac curriculum was found in search results.

\section{Linguistic and Reference Materials}
\begin{itemize}
\item George Kiraz, \emph{Lexical Tools to the Syriac New Testament} (with a skeleton grammar by Sebastian Brock), a Syriac-English dictionary with word frequencies and conjugation tables. \url{https://gorgiaspress.com/lexical-tools-to-the-syriac-new-testament}
\item Gospel Studies listing of Syriac resources, mentioning public-domain grammars (N\"oldeke, Phillips) available online. \url{https://www.biblicalstudies.gospelstudies.org.uk/lang_syriac.php}
\item Ebbe Knudsen, \emph{Classical Syriac Phonology} (Gorgias Press, 2015), mentioned in search results with no direct URL.
\end{itemize}

\section{Religious and Community Resources}
\begin{itemize}
\item Beth Mardutho, the Syriac Institute (Wikipedia entry): publisher of the open-access journal \emph{Hugoye} and sponsor of Syriac computing projects. \url{https://en.wikipedia.org/wiki/Beth_Mardutho}
\end{itemize}

\section{Digital and Audio Resources}
\begin{itemize}
\item Digital Syriac Corpus, project history. \url{https://syriaccorpus.org/history.html}
\item eBethArk\'e: The Syriac Digital Library (Beth Mardutho and Rutgers University Libraries). \url{https://rucore.libraries.rutgers.edu/portals/ebetharke/}
\end{itemize}
No audio or video resources were found.

\section{Curriculum Development Context}
Classical Syriac has solid scholarly and liturgical resources, mostly aimed at university students and in English, and a digital corpus. No primary or secondary school curriculum, or native-language (Syriac-medium) teaching materials, appeared in the search results. A search for Syriac or Suryoyo mother-tongue schooling in Sweden or Iraq returned only university course pages, so that question remains open. A K-12 curriculum would have to decide between Classical Syriac and a modern spoken variety.

\end{document}